\subsection{深度与连通性}

引理 4.--- 设 A 为 Noether 环，F 为 Spec(A) 的闭子集，U 为其补开集，u : M → N 为有限生成 A-模之间的同态。

\begin{enumerate}
\def\labelenumi{\alph{enumi})}
\item
  设 \(u_{\mathfrak{p}}: \mathrm{M}_{\mathfrak{p}} \to \mathrm{N}_{\mathfrak{p}}\) 对所有 \(\mathfrak{p} \in \mathrm{U}\) 单射，并设 \(\operatorname{prof}_{\mathrm{F}}(\mathrm{M}) \geqslant 1\) 。则 \(u\) 单射。
\item
  设 \(u_{\mathfrak{p}}: \mathrm{M}_{\mathfrak{p}} \to \mathrm{N}_{\mathfrak{p}}\) 对所有 \(\mathfrak{p} \in \mathrm{U}\) 双射，并设 \(\operatorname{prof}_{\mathrm{F}}(\mathrm{M}) \geqslant 2\) 与 \(\operatorname{prof}_{\mathrm{F}}(\mathrm{N}) \geqslant 1\) 。则 u 双射。
\item
  诸假设蕴含 Supp(Ker \(u\) ) \(\subset\) F，于是 \(\operatorname{Hom}_{\Lambda}(\operatorname{Ker} u, \mathbf{M}) = 0\) （ \(n^{\circ}\) 5, 注 1）；因此有 \(\operatorname{Ker} u = 0\) 。
\item
  我们已知 \(u\) 单射，且有 \(\operatorname{Supp}(\operatorname{Coker} u) \subset \mathbb{F}\) 。由同上引文，有 \(\operatorname{Hom}_{\mathbf{A}}(\operatorname{Coker} u, \mathbf{N}) = 0\) 与 \(\operatorname{Ext}_{\mathbf{A}}^{1}(\operatorname{Coker} u, \mathbf{M}) = 0\) 。由扩张模的正合序列
\end{enumerate}

\[
\operatorname{Hom} _ {\mathrm{A}} (\operatorname{Coker} u, \mathrm{N}) \rightarrow \operatorname{Hom} _ {\mathrm{A}} (\operatorname{Coker} u, \operatorname{Coker} u) \rightarrow \operatorname{Ext} _ {\mathrm{A}} ^ {1} (\operatorname{Coker} u, \mathrm{M})
\]

我们推出 \(\operatorname{Hom}_{\mathrm{A}}(\operatorname{Coker} u, \operatorname{Coker} u) = 0\) ，这蕴含 \(\operatorname{Coker} u = 0\) 。

注 1.—设 A 为 Noether 环，F 为 \(\operatorname{Spec}(A)\) 的闭子集，U 为其补开集。要使 \(\operatorname{prof}_{\mathrm{F}}(A) \geqslant 1\) ，必须且只需 \(\operatorname{Ass}(A) \subset U\) （注 2, 第 5 节）。当此条件满足时，\(\operatorname{Spec}(A)\) 的每个不可约分支都与 U 相交，故 U 在 \(\operatorname{Spec}(A)\) 中稠密。

定理 3 (Hartshorne).--- 设 A 为 Noether 环，F 为 Spec(A) 的闭子集，U 为其补开集。设 prof \(_{F}\) (A) ≥ 2。则对 Spec(A) 的每个连通分支 Y，集合 Y ∩ U 连通且在 Y 中稠密。

首先设 \(\operatorname{Spec}(A)\) 连通。由注 1，\(U\) 在 \(\operatorname{Spec}(A)\) 中稠密，剩下要证明它连通。反设两个不相交的非空开子集 \(U_0\) 与 \(U_1\) —— \(\operatorname{Spec}(A)\) 的开子集——之并为 \(U\)。由注 1，\(\operatorname{Ass}(A)\) 包含于 \(U\)，故它是 \(\operatorname{Ass}(A) \cap U_0\) 与 \(\operatorname{Ass}(A) \cap U_1\) 的不交并。由 IV, § 1, 第 1 节, 命题 4，存在 A 的理想 \(J_0\) 与 \(J_1\) 使 \(\operatorname{Ass}(J_i) = \operatorname{Ass}(A) \cap U_i\) ， \(\operatorname{Ass}(A/J_i) = \operatorname{Ass}(A) \cap U_{1-i}\quad (i=0,1)\) 。\(U_i\) 在 \(\operatorname{Spec}(A)\) 中的补集包含 \(\operatorname{Ass}(A/J_i)\) 与 \(\operatorname{Ass}(J_{1-i})\)；由于它是闭集，它包含 \(\operatorname{Supp}(A/J_i)\) 与 \(\operatorname{Supp}(J_{1-i})\)。于是对 \(\mathfrak{p} \in U_i\) ，有 \((J_i)_{\mathfrak{p}} = A_{\mathfrak{p}}\) 与 \((J_{1-i})_{\mathfrak{p}} = 0\) ；特别地这蕴含 \(J_0\) 与 \(J_1\) 异于 A。设 B 为 A-模 \(A/J_0 \times A/J_1\) ，设 \(u: A \to B\) 为典范同态。由前所述，同态 \(u_{\mathfrak{p}}\) 对每个 \(\mathfrak{p} \in U\) 双射；此外有 \(\operatorname{Ass}(B) \subset U_0 \cup U_1 = U\) ，故由注 1 有 \(\operatorname{prof}_F(B) \geqslant 1\) 。于是引理 4 蕴含 \(u\) 双射，这与 \(\operatorname{Spec}(A)\) 的连通性矛盾。

现在处理一般情形。设 J 为 A 的理想使 \(F = V(J)\)，并设 Y 为 \(\operatorname{Spec}(A)\) 的一个连通分支。由 II, § 4, 第 3 节, 命题 15，存在 A 的幂等元 \(f\) 使 Y 等同于 \(\operatorname{Spec}(A_f)\) （ \(\operatorname{Spec}(A)\) 的一个部分）。于是 \(Y \cap F\) 等同于 \(V(J_f)\)；由第 3 节的命题 6 a)，有 \(\operatorname{prof}_{A_f}(J_f, A_f) \geqslant \operatorname{prof}_A(J; A) \geqslant 2\) 。由证明的第一部分得出，\(Y \cap U = Y - (Y \cap F)\) 连通且在 Y 中稠密。

推论 1.--- 把 U 的每个连通分支映为其在 Spec(A) 中的闭包的映射，是从 U 的连通分支组成的集合到 Spec(A) 的连通分支组成的集合上的一个双射。

推论 2.--- 对每个深度 \(\geqslant\) 2 的 Noether 局部环 B，拓扑空间 \(\operatorname{Spec}(\mathbf{B}) = \{\mathfrak{m}_{\mathbf{B}}\}\) 连通。

推论 3.--- 在定理 3 的假设下，设 \(\text{Spec}(A_{\mathfrak{p}})\) 不可约（相应地，\(A_{p}\) 为整环）对一切 \(p \in U\) 成立；则 \(\text{Spec}(A_{\mathfrak{p}})\) 不可约（相应地，\(A_{p}\) 为整环）对一切 \(p \in \text{Spec}(A)\) 成立。

设 \((\mathrm{Y}_{i})_{i\in\mathrm{I}}\) 为 \(\operatorname{Spec}(A)\) 的不可约分支组成的（有限）族。设 \(\mathfrak{p}\in U\) ；由于 \(\operatorname{Spec}(A_{\mathfrak{p}})\) 不可约，\(\mathfrak{p}\) 只包含 A 的一个极小素理想，因而只属于诸 \(Y_i\) 之一 (II, § 4, 第 3 节, 命题 14 的推论 2)。子集 \(Y_i\cap U\) 是 U 的闭子集，两两不交，且因 U 在 \(\operatorname{Spec}(A)\) 中稠密而非空，并由 II, § 4, 第 1 节, 命题 7 不可约；因而它们是 U 的连通分支。它们的闭包 \(Y_i\) 是 \(\operatorname{Spec}(A)\) 的连通分支（推论 1）。这就证明了 \(\operatorname{Spec}(A)\) 的连通分支不可约，从而 \(\operatorname{Spec}(A_{\mathfrak{p}})\) 对每个 \(\mathfrak{p}\) 不可约（第 8 节）。

设 \(A_{\mathfrak{q}}\) 对所有 \(\mathfrak{q} \in U\) 是整的。设 \(\mathfrak{p} \in \operatorname{Spec}(A)\) 。由于 \(\operatorname{Spec}(A_{\mathfrak{p}})\) 不可约，\(A_{\mathfrak{p}}\) 的诣零根是 \(A_{\mathfrak{p}}\) 的唯一极小素理想；因此它属于 \(\operatorname{Ass}(A_{\mathfrak{p}})\) (IV, § 1, 第 3 节, 命题 7 的推论 1)，从而等于 \(\mathfrak{q}A_{\mathfrak{p}}\) ，其中 \(\mathfrak{q}\) 是与 A 相伴的素理想 (IV, § 1, 第 2 节, 命题 5 的推论)。有 \(\mathfrak{q} \in U\) （注 1）且 \(\mathfrak{q}A_{\mathfrak{q}} \in \operatorname{Ass}(A_{\mathfrak{q}})\) （同上引文）；由于 \(A_{\mathfrak{q}}\) 是整的，\(\mathfrak{q}\) 为零，因此 \(A_{\mathfrak{p}}\) 是整环。

推论 4.--- 设 \(A\) 为谱连通的 Noether 环。设存在整数 \(d \geqslant 1\) ，使 \(\operatorname{prof}(A_{\mathfrak{p}}) \geqslant 2\) 对 A 的每个素理想 \(\mathfrak{p}\) （高度 \textgreater{} d 者）成立。

\begin{enumerate}
\def\labelenumi{\alph{enumi})}
\item
  对 \(\operatorname{Spec}(\mathbf{A})\) 的每个余维数 \(>d\) 的闭子集 Z，空间 \(\operatorname{Spec}(\mathbf{A}) - \mathbf{Z}\) 连通。
\item
  设 Y 与 Y' 为 Spec(A) 的不可约分支。存在 Spec(A) 的不可约分支的序列 \(X_{1}, X_{2}, \ldots, X_{n}\) ，使得有 \(X_{1} = Y\) ， \(X_{n} = Y'\) ，且对 \(i = 1, \ldots, n - 1\)
\end{enumerate}

\[
\operatorname{codim} \left(\mathrm{X} _ {i} \cap \mathrm{X} _ {i + 1}, \operatorname{Spec} (\mathrm{A})\right) \leqslant d.
\]

设 \(Z \subset \operatorname{Spec}(A)\) 为余维数 \(>d\) 的闭子集。对每个 \(\mathfrak{p} \in Z\) ，有 \(\dim(A_{\mathfrak{p}}) > d\) ，从而 \(\operatorname{prof}(A_{\mathfrak{p}}) \geqslant 2\) ，这蕴含 \(\operatorname{prof}_{Z}(A)\) 为 \(\geqslant 2\)（第 5 节, 命题 8），从而 \(\operatorname{Spec}(A) - Z\) 连通（定理 3）。

设 Z 为诸集合 \(X^{\prime} \cap X^{\prime\prime}\) 之并，其中 \((X^{\prime}, X^{\prime\prime})\) 遍历 Spec(A) 的不可约分支的（有限）对组成的集合中满足 \(\text{codim}(X^{\prime} \cap X^{\prime\prime}, \text{Spec}(A)) > d\) 者。由 a)，集合 \(\text{Spec}(A) - Z\) 连通。\(\text{Spec}(A)\) 的所有不可约分支都与 \(\text{Spec}(A) - Z\) 相交；它们在 \(\text{Spec}(A) - Z\) 上的迹是 \(\text{Spec}(A) - Z\) 的不可约分支 (II, § 4, 第 1 节, 命题 7)。由于 \(\text{Spec}(A) - Z\) 连通，存在序列 \(X_{1}, \ldots, X_{n}\) ，由 \(\text{Spec}(A)\) 的不可约分支组成，使得有 \(X_{1} - Z = Y - Z\) ， \(X_{n} - Z = Y' - Z\) ，且 \((X_{i} - Z) \cap (X_{i+1} - Z) \neq \emptyset\) 对 \(1 \leqslant i \leqslant n - 1\) 成立；换言之，有 \(X_{1} = Y\) ， \(X_{n} = Y'\) 且 \(\text{codim}(X_{i} \cap X_{i+1}) \leqslant d\) 。

注 2. 当 A 是 Macaulay 环时，可对 \(d=1\) 应用该推论（ \(\S\) 2, 第 5 节）。

例（由 4 维仿射空间的四个坐标平面构成的完全交）.--- 设 \(k\) 为域。用 S 表示多项式环 \(k[\mathrm{T}_{1},\mathrm{T}_{2},\mathrm{T}_{3},\mathrm{T}_{4}]\)。回忆 (VIII, § 2, 第 4 节, 定理 3)，S 的素理想的每个极大链长度为 4。设 \(\mathfrak{m}\) 为 S 的由诸 \(\mathrm{T}_{i}\) 生成的理想，\(\mathfrak{a}\) 为由 \(\mathrm{T}_{1}\mathrm{T}_{2}\) 与 \(\mathrm{T}_{3}\mathrm{T}_{4}\) 生成的理想，且对 \(1 \leqslant i < j \leqslant 4\) ，设 \(\mathfrak{p}_{ij}\) 为由 \(\mathrm{T}_{i}\) 与 \(\mathrm{T}_{j}\) 生成的理想。诸理想 \(\mathfrak{p}_{ij}\) 是高度 2 的素理想，它们的和是极大理想 \(\mathfrak{m}\) ，且有 \(\mathfrak{a} = \mathfrak{p}_{13} \cap \mathfrak{p}_{14} \cap \mathfrak{p}_{23} \cap \mathfrak{p}_{24}\)。

\begin{enumerate}
\def\labelenumi{\alph{enumi})}
\item
  环 A = S/a 是约化的；Spec(A) 的不可约分支是诸集合 \(X_{ij} = V(\mathfrak{p}_{ij}/\mathfrak{a})\) （i = 1, 2, j = 3, 4），它们维数为 2 且都包含点 m/a。特别地，Spec(A) 连通且维数为 2。两个不同分支 \(X_{ij}\) 与 \(X_{kl}\) 的交退化为 \(\{m/a\}\) 当 \(\{i,j\} \cap \{k,l\} = \varnothing\) 时，否则维数为 1。由此得出，推论 4 的结论对 d = 1 成立（我们后面将看到 A 是 Macaulay 环，因而推论 4 的假设本身对 d = 1 成立）。
\item
  设 b 为 S 的由 \(T_{1}T_{2}\) 、 \(T_{1}T_{3}\) 、 \(T_{2}T_{4}\) 、 \(T_{3}T_{4}\) 生成的理想，B 为环 S/b。有 \(b = p_{14}p_{23} = p_{14} \cap p_{23}\) 。环 B 是约化的。它的谱等同于 Spec(A) 的闭子集 \(X_{14} \cup X_{23}\) ；它有两个不可约分支（维数为 2），其交退化为一点。B 沿这一点的深度因 B 约化而严格为正，且由定理 3 知小于 2，故等于 1。于是 B 不是 Macaulay 环。
\end{enumerate}

